% JAIR reproducibility checklist. The item text is the AuthorKit's, copied
% unmodified so a reviewer can diff it against the template; only the answers
% and the short justifications are ours.
%
% This is a survey, so two of the four blocks do not apply, and saying so is
% part of the answer. The paper runs no algorithm and reports no measured
% quantity of its own, so the computational-experiments block is answered "no"
% at the gate. It does introduce a data set, the rubric score file, so the
% data-set block is answered in full. Do not soften a "no" here into a
% "partially" to make the appendix look better: the gate questions are the
% ones a reader uses to calibrate the rest.

% The heading names JAIR in the submission build, because the checklist is the
% journal's own instrument and the AuthorKit's wording. The arXiv build drops
% the venue name: the checklist is still worth publishing, and the preprint
% should not assert where it was submitted. \ifarxiv is set in main.tex.
\ifarxiv
  \section{Reproducibility Checklist}
\else
  \section{Reproducibility Checklist for JAIR}
\fi
\label{sec:appendix:repro}

Select the answers that apply to your research, one per item.

\subsection*{All articles:}

\begin{enumerate}
    \item All claims investigated in this work are clearly stated.
    [\textbf{yes}]
    \item Clear explanations are given how the work reported substantiates the claims.
    [\textbf{yes}]
    \item Limitations or technical assumptions are stated clearly and explicitly.
    [\textbf{yes}]
    \item Conceptual outlines and/or pseudo-code descriptions of the AI methods introduced in this work are provided, and important implementation details are discussed.
    [\textbf{yes}]
    \item
    Motivation is provided for all design choices, including algorithms, implementation choices, parameters, data sets and experimental protocols beyond metrics.
    [\textbf{yes}]
\end{enumerate}

\noindent The rubric's ten axes, their ordinal levels, and the scoring
procedure are stated in \S\ref{sec:rubric}, and each axis level is anchored to
a specification property rather than to a judgement of quality. The
limitations are gathered in \S\ref{sec:rubric:limits} and are not confined
there: the single-rater design and the absence of an inter-rater agreement
statistic are stated in the abstract, in \S\ref{sec:rubric:calibration}, and
again in the conclusion. The three proposed extensions of
\S\ref{sec:proposals} are specification sketches, not validated standards,
and are labelled as such wherever they are introduced.

\subsection*{Articles containing theoretical contributions:}
Does this paper make theoretical contributions?
[\textbf{no}]

\noindent The survey proves nothing. The three proposed extensions of
\S\ref{sec:proposals} state observable invariants that a receiver can check,
but those are specification requirements, not theorems, and the paper does
not claim them as results.

\subsection*{Articles reporting on computational experiments:}
Does this paper include computational experiments? [\textbf{no}]

\noindent No algorithm is run and no performance quantity is measured. Every
number in the paper is either a rubric score assigned by a rater reading a
specification (\S\ref{sec:rubric:method}), a count over those scores, or a
result quoted from the cited source with attribution. The quoted empirical
results, such as the leakage rates in \S\ref{sec:security}, are other
authors' measurements and are reproducible only against their artifacts, not
ours.

\subsection*{Articles using data sets:}
Does this work rely on one or more data sets (possibly obtained from a
benchmark generator or similar software artifact)?
[\textbf{yes}]

\begin{enumerate}
    \item
    All newly introduced data sets
    are included in an online appendix
    or will be made publicly available upon publication of the paper.
    The online appendix follows best practices for long-term accessibility with a license
    that allows free usage for research purposes.
    [\textbf{yes}]
    \item The newly introduced data set comes with a license that
    allows free usage for reproducibility purposes.
    [\textbf{yes}]
    \item The newly introduced data set comes with a license that
    allows free usage for research purposes in general.
    [\textbf{yes}]
    \item All data sets drawn from the literature or other public sources (potentially including authors' own previously published work) are accompanied by appropriate citations.
    [\textbf{yes}]
    \item All data sets drawn from the existing literature (potentially including authors' own previously published work) are publicly available. [\textbf{yes}]
    \item All new data sets and data sets that are not publicly available are described in detail, including relevant statistics, the data collection process and annotation process if relevant.
    [\textbf{yes}]
    \item
    All methods used for preprocessing, augmenting, batching or splitting data sets (e.g., in the context of hold-out or cross-validation)
    are described in detail. [\textbf{NA}]
\end{enumerate}

\noindent The newly introduced data set is the rubric itself as data:
\texttt{protocol-scores.yaml} (one entry per protocol and axis, each carrying
the score, a one-sentence evidence summary, and the citation the score rests
on), \texttt{threat-coverage.yaml}, and \texttt{industry-adoption.yaml}. It is
Contribution C4 (\S\ref{sec:intro:contributions}). Its collection procedure is
the scoring procedure of \S\ref{sec:rubric:method}, which is the annotation
protocol for this data set and is stated as a numbered sequence there. The
companion repository also carries the scripts that turn the three files into
Table~\ref{tab:rubric-matrix}, Table~\ref{tab:threat-surface},
Table~\ref{tab:industry-adoption}, Table~\ref{tab:severity-coverage},
Figure~\ref{fig:heatmap}, Figure~\ref{fig:radar}, and
Figure~\ref{fig:taxonomy}, so every generated artifact in the paper is
re-derivable from the data. The repository is at \companionrepo.

The last item is \textbf{NA} because there is no split. The data set is a
census, not a sample: every axis of every scored protocol is populated, so
there is nothing held out and nothing to cross-validate. Which protocols
enter the set, and why the thirteen are these thirteen, is stated in
\S\ref{sec:taxonomy} rather than left to the reader.

\subsection*{Explanations on any of the answers above (optional):}

Two answers need more than a yes or no.

On the data-set license: the three data files are released under the
Creative Commons Attribution 4.0 license, and the scripts under the MIT
License. Both permit free use for research and for reproducibility. Credit is
the only condition. The JAIR class files in the repository keep their
original terms.

On coverage over time: the survey declares a 2026-Q1 cutoff on page 1, and
protocol specifications move faster than a review cycle. Two protocols in the
scored set, IBM ACP and AGNTCY/SLIM, were scored in a later pass than the
other six and are dated separately in \texttt{protocol-scores.yaml}
(\S\ref{sec:comparison:matrix}). Reproducing a score means reading the
specification version the entry names, not the current one.
